\section{Scalar spectral jets and ordinary resolvent integrals}
\label{sp:section}

Throughout this section $a>0$, $x_n=n+a$ for $n=0,1,\ldots$, and
$\mu_1,\ldots,\mu_N$ are positive real numbers, with $N\geq1$.
The logarithm $\Log$ is the principal logarithm on
$\mathbb C\setminus(-\infty,0]$.  Powers of positive real numbers use
their real logarithms.  We write $\zeta(s,a)$ for the Hurwitz zeta
function, $\psi$ for the digamma function, and
$\zeta^{(r)}(s,a)=\partial_s^r\zeta(s,a)$.
The formulas below are specialized deductions for scalar factors and
their connected logarithmic products.  They make no assertion of
global priority for the general principle relating spectral finite
parts and local asymptotic data.

\subsection{The scalar master formula}
\label{sp:scalar-subsection}

For $\Re s>1/(2N)$ define
\begin{equation}
 Z_N(s;a,\boldsymbol\mu)
 :=\sum_{n=0}^{\infty}\prod_{j=1}^{N}(x_n^2+\mu_j)^{-s}.
 \label{sp:zeta-definition}
\end{equation}
For integers $r\geq1$, its $r$th spectral jet is
$J_{N,r}(a;\boldsymbol\mu):=\partial_s^rZ_N(0;a,\boldsymbol\mu)$,
once continuation to a neighborhood of zero has been established.

\begin{lemma}[Continuation by a normally convergent difference]
\label{sp:continuation}
The function in \eqref{sp:zeta-definition} has the continuation
\begin{equation}
 Z_N(s;a,\boldsymbol\mu)
 =\zeta(2Ns,a)+
 \sum_{n=0}^{\infty}x_n^{-2Ns}
 \left\{\exp\left[-s\sum_{j=1}^{N}
             \log\left(1+\frac{\mu_j}{x_n^2}\right)\right]-1\right\},
 \label{sp:normal-continuation}
\end{equation}
where the series is normally convergent on $\Re s>-1/(2N)$.
In particular $Z_N$ is holomorphic near $s=0$, $Z_N(0)=\zeta(0,a)$,
and
\begin{equation}
 J_{N,r}=(2N)^r\zeta^{(r)}(0,a)
          +(-1)^r\sum_{n=0}^{\infty}F_{N,r}(x_n^2),
 \qquad
 F_{N,r}(z):=
 \left(\sum_{j=1}^{N}\Log(z+\mu_j)\right)^r-N^r\Log^r z.
 \label{sp:jet-ordinary-series}
\end{equation}
The series in \eqref{sp:jet-ordinary-series} is absolutely convergent.
\end{lemma}

\begin{proof}
Put $L_n=\sum_j\log(1+\mu_j/x_n^2)$.  Then $L_n=O(x_n^{-2})$.
For a compact subset $K$ of $\Re s>-1/(2N)$, the tail of the
series in \eqref{sp:normal-continuation} is bounded, uniformly on $K$,
by a constant times $x_n^{-2N\inf_K\Re s-2}$.  This is summable.
The finitely many remaining summands are entire functions of $s$.
The identity first holds in the defining half-plane by factoring
out $x_n^{-2Ns}$, and hence gives the asserted continuation.
The sole pole of $\zeta(2Ns,a)$ in this region is away from zero.
Normal convergence permits differentiation at zero, giving
\eqref{sp:jet-ordinary-series}.  Finally,
\[
 \sum_j\log(x_n^2+\mu_j)=N\log(x_n^2)+O(x_n^{-2}),
 \qquad
 F_{N,r}(x_n^2)=O\left(\frac{1+\log^{r-1}x_n}{x_n^2}\right)
\]
for large $n$, which also proves absolute convergence directly.
\end{proof}

For $t>0$ distinct from the masses, set
\begin{equation}
 A(t):=\sum_{j=1}^{N}\log|\mu_j-t|,
 \qquad k(t):=\#\{j:\mu_j<t\},
 \label{sp:cut-data}
\end{equation}
counting repeated masses with their multiplicity, and define the
real, possibly signed density
\begin{equation}
 \rho_{N,r}(t):=-\frac{1}{\pi}\Im\left\{
        [A(t)+i\pi k(t)]^r
           -N^r[\log t+i\pi]^r\right\}.
 \label{sp:scalar-density}
\end{equation}
Its values at the finitely many mass points may be assigned
arbitrarily.  They do not change any integral below.

\begin{theorem}[A single ordinary integral for every scalar jet]
\label{sp:master}
For every integer $N\geq1$, every integer $r\geq1$, every $a>0$,
and all positive masses,
\begin{equation}
 \boxed{\displaystyle
 J_{N,r}(a;\boldsymbol\mu)
 =(2N)^r\zeta^{(r)}(0,a)
       +(-1)^r\int_0^{\infty}\rho_{N,r}(t)Q_a(t)\,dt,}
 \label{sp:master-formula}
\end{equation}
where
\begin{equation}
 Q_a(t):=\sum_{n=0}^{\infty}\frac{1}{x_n^2+t}
 =\frac{\psi(a+i\sqrt t)-\psi(a-i\sqrt t)}{2i\sqrt t},
 \qquad Q_a(0):=\psi'(a).
 \label{sp:resolvent-kernel}
\end{equation}
The integral is absolutely convergent.  At the mass points it is
understood as the sum of the ordinary improper integrals over the
intervals separated by the distinct masses; no finite part or
principal value is required.
\end{theorem}

\begin{proof}
We first establish the Stieltjes representation
\begin{equation}
 F_{N,r}(z)=\int_0^{\infty}\frac{\rho_{N,r}(t)}{z+t}\,dt,
 \qquad z\in\mathbb C\setminus(-\infty,0].
 \label{sp:scalar-stieltjes}
\end{equation}
The function $F_{N,r}$ is holomorphic on this slit plane and
satisfies Schwarz reflection.  On the upper bank of the cut,
\[
 \sum_j\Log(\mu_j-t+i0)=A(t)+i\pi k(t),
 \qquad \Log(-t+i0)=\log t+i\pi.
\]
Consequently its upper boundary value has imaginary part
$-\pi\rho_{N,r}(t)$.  Near zero,
$F_{N,r}(z)=O(1+|\Log z|^r)$; at a nonzero endpoint $-\mu_j$
it has at worst a power of a logarithm.  The integrals on circles
of radius $\varepsilon$ around these endpoints therefore tend to
zero, since $\varepsilon(1+|\log\varepsilon|^r)\to0$.
At infinity, uniformly along a large keyhole contour,
\[
 \sum_j\Log(z+\mu_j)
       =N\Log z+O(|z|^{-1}),
 \qquad
 F_{N,r}(z)=O\left(\frac{1+\log^{r-1}|z|}{|z|}\right).
\]
The outer circle in the Cauchy integral hence contributes zero
in the limit.  Cauchy's formula, followed by collapsing the two
banks of the cut, gives \eqref{sp:scalar-stieltjes} with the sign
in \eqref{sp:scalar-density}.

Here are explicit bounds that justify the resulting improper
integral and its later interchange with a series.  Since the
masses are positive, $k(t)=0$ for sufficiently small $t$ and
$A(t)$ remains bounded there.  Thus
\begin{equation}
 \rho_{N,r}(t)=O(1+|\log t|^{r-1})\qquad(t\downarrow0).
 \label{sp:density-zero-bound}
\end{equation}
At each mass point the density is bounded by a constant times
$1+|\log|t-\mu_j||^{r-1}$, allowing for multiplicities; these
singularities are locally integrable.  Beyond the largest mass,
write
\[
 A(t)=N\log t+\delta(t),\qquad
 \delta(t)=\sum_j\log(1-\mu_j/t)=O(t^{-1}),\qquad k(t)=N.
\]
In the difference
$[N(\log t+i\pi)+\delta(t)]^r-[N(\log t+i\pi)]^r$,
taking the imaginary part lowers the maximal logarithmic degree
by one.  Hence, for $r\geq2$,
\begin{equation}
 \rho_{N,r}(t)=O\left(\frac{1+\log^{r-2}t}{t}\right)
       \qquad(t\to\infty).
 \label{sp:density-infinity-bound}
\end{equation}
For $r=1$ the density is zero beyond the largest mass.
These estimates also directly give absolute convergence in
\eqref{sp:scalar-stieltjes} at any fixed $z$ off the cut.

The convergent digamma difference series gives
\eqref{sp:resolvent-kernel}.  It also shows that $Q_a(t)>0$,
that $Q_a$ is continuous at zero with value $\psi'(a)$, and that
it is bounded on compact intervals of $t\geq0$.
Comparison of its decreasing summand with an integral yields
\[
 Q_a(t)=\int_a^{\infty}\frac{dx}{x^2+t}+O(t^{-1})
       =\frac{\pi}{2\sqrt t}+O(t^{-1})
       \qquad(t\to\infty).
\]
It follows from the endpoint bounds that
\[
 \int_0^{\infty}|\rho_{N,r}(t)|Q_a(t)\,dt<\infty.
\]
Tonelli's theorem applied to the absolute values, followed by
Fubini's theorem, now proves
\[
 \sum_{n=0}^{\infty}F_{N,r}(x_n^2)
   =\int_0^{\infty}\rho_{N,r}(t)
                     \sum_{n=0}^{\infty}\frac{dt}{x_n^2+t}
   =\int_0^{\infty}\rho_{N,r}(t)Q_a(t)\,dt.
\]
Substitution in \eqref{sp:jet-ordinary-series} gives the result.
\end{proof}

\begin{remark}[Explicit densities and collisions]
\label{sp:explicit-densities}
Expanding only the odd binomial terms gives, at every order,
\begin{equation}
 \rho_{N,r}(t)=
 \sum_{h=0}^{\lfloor(r-1)/2\rfloor}
   (-1)^h\binom{r}{2h+1}\pi^{2h}
   \left\{N^r(\log t)^{r-2h-1}
                -k(t)^{2h+1}A(t)^{r-2h-1}\right\}.
 \label{sp:binomial-density}
\end{equation}
In particular,
\begin{align}
 \rho_{N,1}(t)&=N-k(t),\label{sp:density-one}\\
 \rho_{N,2}(t)&=2N^2\log t-2k(t)A(t),\label{sp:density-two}\\
 \rho_{N,3}(t)&=3N^3\log^2t-3k(t)A(t)^2
                    +\pi^2\{k(t)^3-N^3\}.
 \label{sp:density-three}
\end{align}
The densities need not be nonnegative at higher orders.
Repeated masses are included directly in these formulas: $k$
jumps by their multiplicity, and $A$ contains the corresponding
multiple of $\log|\mu-t|$.  The proof of
Theorem~\ref{sp:master} applies without a distinctness assumption,
so no limiting prescription is needed when masses coincide.
\end{remark}

\begin{corollary}[The first jet and its normalization]
\label{sp:first-jet}
Let $\log\Gamma$ denote the analytic branch on $\Re z>0$ that
is real on the positive real axis.  Then
\begin{equation}
 J_{N,1}(a;\boldsymbol\mu)
 =\sum_{j=1}^{N}\left\{
       \log\Gamma(a+i\sqrt{\mu_j})
       +\log\Gamma(a-i\sqrt{\mu_j})-\log(2\pi)\right\}.
 \label{sp:first-jet-gamma}
\end{equation}
\end{corollary}

\begin{proof}
Since $N-k(t)=\sum_j\boldsymbol1_{(0,\mu_j)}(t)$, the integral
in \eqref{sp:master-formula} is a sum of $\int_0^{\mu_j}Q_a(t)\,dt$.
Differentiation with respect to $t$ gives
\[
 \frac{d}{dt}\left\{
   \log\Gamma(a+i\sqrt t)+\log\Gamma(a-i\sqrt t)\right\}
 =-Q_a(t).
\]
The identity extends continuously to $t=0$.  Therefore
\[
 \int_0^\mu Q_a(t)\,dt
 =2\log\Gamma(a)-\log\Gamma(a+i\sqrt\mu)
                       -\log\Gamma(a-i\sqrt\mu).
\]
Combining this with
$\zeta'(0,a)=\log\Gamma(a)-\tfrac12\log(2\pi)$ proves
\eqref{sp:first-jet-gamma}.
\end{proof}

\subsection{Parameter calculus without regularized interchange}
\label{sp:parameters}

\begin{proposition}[Derivatives and an anchored primitive in $a$]
\label{sp:a-calculus}
For every nonnegative integer $m$,
\begin{equation}
 \partial_a^mJ_{N,r}
 =(2N)^r\partial_a^m\zeta^{(r)}(0,a)
   +(-1)^r\int_0^{\infty}\rho_{N,r}(t)Q_a^{[m]}(t)\,dt,
 \label{sp:a-derivatives}
\end{equation}
where, with $\psi^{(0)}=\psi$,
\begin{equation}
 Q_a^{[m]}(t):=
 \frac{\psi^{(m)}(a+i\sqrt t)-\psi^{(m)}(a-i\sqrt t)}
      {2i\sqrt t},
 \qquad Q_a^{[m]}(0)=\psi^{(m+1)}(a).
 \label{sp:polygamma-kernel}
\end{equation}
All these integrals are absolutely convergent and locally uniform
for $a>0$.

For a fixed anchor $a_*>0$, define
\begin{align}
 K_a(t)&:=\frac{\log\Gamma(a+i\sqrt t)-\log\Gamma(a-i\sqrt t)}
                      {2i\sqrt t},\qquad t>0,\notag\\
 H_{a,a_*}(t)&:=K_a(t)-K_{a_*}(t),
 \label{sp:gamma-kernel}\\
 H_{a,a_*}(0)&:=\psi(a)-\psi(a_*),\notag\\
 B_{N,r}(a)
 &:=(2N)^r r!\sum_{j=0}^{r}\frac{\zeta^{(j)}(-1,a)}{j!}.
 \label{sp:primitive-baseline}
\end{align}
The branch of $\log\Gamma$ is the one used in
Corollary~\ref{sp:first-jet}.  Then the anchored primitive is
\begin{equation}
 \int_{a_*}^{a}J_{N,r}(u;\boldsymbol\mu)\,du
 =B_{N,r}(a)-B_{N,r}(a_*)
     +(-1)^r\int_0^{\infty}
                      \rho_{N,r}(t)H_{a,a_*}(t)\,dt.
 \label{sp:anchored-primitive}
\end{equation}
The final integral is again ordinary and absolutely convergent.
\end{proposition}

\begin{proof}
Let $a$ range over a compact interval $I\subset(0,\infty)$.
Choose a positive complex radius $\eta$ smaller than one half of
$\inf I$.  For real $a\in I$, $t\geq0$, and $|w-a|\leq\eta$,
each factor $n+w\pm i\sqrt t$ has modulus at least a fixed
positive multiple of $\sqrt{(n+a)^2+t}$.  Cauchy's derivative
estimate applied to $((n+w)^2+t)^{-1}$ therefore gives
\[
 \left|\partial_a^m\frac{1}{(n+a)^2+t}\right|
 \leq\frac{C_{I,m}}{(n+a)^2+t}.
\]
Summing and using the bounds in the proof of
Theorem~\ref{sp:master} produces a majorant for every derivative
that is integrable against $|\rho_{N,r}(t)|\,dt$, uniformly for
$a\in I$.  Differentiation under the integral is valid to every
order.  Differentiating the digamma expression for $Q_a$ gives
\eqref{sp:polygamma-kernel}; its value at zero follows by a Taylor
expansion.  The same local uniform estimates justify derivatives
of the normally convergent series in
\eqref{sp:normal-continuation}.

For the primitive, integration of $Q_u(t)$ in $u$ gives
\[
 H_{a,a_*}(t)=\int_{a_*}^{a}Q_u(t)\,du.
\]
This identity both fixes the branch normalization and proves the
required integral bound, including the continuous value at $t=0$.
For the Hurwitz term, the identity
$\partial_a\zeta(v,a)=-v\zeta(v+1,a)$ gives
\[
 \partial_a\left\{-\frac{\zeta(2Ns-1,a)}{2Ns-1}\right\}
 =\zeta(2Ns,a).
\]
Writing $u=2Ns$ and expanding at $u=0$,
\[
 -\frac{\zeta(2Ns-1,a)}{2Ns-1}
 =\frac{\zeta(u-1,a)}{1-u},\qquad
 \left.\partial_s^r
    \frac{\zeta(2Ns-1,a)}{1-2Ns}\right|_{s=0}
 =(2N)^r r!\sum_{j=0}^{r}\frac{\zeta^{(j)}(-1,a)}{j!}.
\]
Thus $B_{N,r}'(a)=(2N)^r\zeta^{(r)}(0,a)$.
Integrating \eqref{sp:master-formula}, with the locally uniform
absolute bounds just proved, gives
\eqref{sp:anchored-primitive}.
\end{proof}

\begin{proposition}[Mixed mass derivatives and resolvents]
\label{sp:mass-derivatives}
Hold $N$ fixed.  Choose $q$ distinct labels
$i_1,\ldots,i_q\in\{1,\ldots,N\}$ and put
$S_n:=\sum_{j=1}^{N}\log(x_n^2+\mu_j)$.
For $1\leq q\leq r$,
\begin{equation}
 \partial_{\mu_{i_1}}\cdots\partial_{\mu_{i_q}}J_{N,r}
 =(-1)^r\frac{r!}{(r-q)!}
   \sum_{n=0}^{\infty}
     \frac{S_n^{r-q}}{\prod_{\ell=1}^{q}(x_n^2+\mu_{i_\ell})}.
 \label{sp:mixed-mass-formula}
\end{equation}
For $q>r$ the mixed derivative vanishes.  Every series in
\eqref{sp:mixed-mass-formula} is absolutely convergent, including
when different labels carry equal masses.
In particular, if $r\leq N$, then
\begin{equation}
 \partial_{\mu_{i_1}}\cdots\partial_{\mu_{i_r}}J_{N,r}
 =(-1)^r r!\sum_{n=0}^{\infty}
          \frac{1}{\prod_{\ell=1}^{r}(x_n^2+\mu_{i_\ell})}
 =-r!\,Q_a[\mu_{i_1},\ldots,\mu_{i_r}],
 \label{sp:rational-mass-resolvent}
\end{equation}
where brackets denote the ordinary divided difference of order
$r-1$ with its confluent continuation.
\end{proposition}

\begin{proof}
In a common domain of absolute convergence, differentiating the
defining spectral summand in $q$ distinct mass labels gives
\[
 \partial_{\mu_{i_1}}\cdots\partial_{\mu_{i_q}}Z_N(s)
 =(-s)^q\sum_{n=0}^{\infty}
     \frac{e^{-sS_n}}{\prod_{\ell=1}^{q}(x_n^2+\mu_{i_\ell})}.
\]
For $q\geq1$ the series on the right is normally convergent in
a neighborhood of $s=0$, uniformly for masses in compact subsets
of the positive orthant.  It supplies the continuation of the
mass derivative.  Taking its $r$th derivative at zero proves
\eqref{sp:mixed-mass-formula} and the vanishing for $q>r$.
The summands are bounded by a constant times
$x_n^{-2q}(1+\log^{r-q}x_n)$ when $q\leq r$.
Finally,
\[
 \left[t\longmapsto\frac{1}{x^2+t}\right]
        [\nu_1,\ldots,\nu_r]
 =\frac{(-1)^{r-1}}{\prod_{j=1}^{r}(x^2+\nu_j)}.
\]
Summing this identity proves the divided-difference formula.
Normal convergence in the masses justifies confluence at equal
arguments.
\end{proof}

\begin{remark}[Repeated differentiation of one mass]
\label{sp:repeated-mass-derivatives}
The qualification ``distinct labels'' in
Proposition~\ref{sp:mass-derivatives} concerns the differentiation
variables, not their numerical values.  A formula for an arbitrary
nonzero multi-index $\boldsymbol\alpha=(\alpha_1,\ldots,\alpha_N)$
is obtained without differentiating a moving branch-cut density.
Let $d=\sum_j\alpha_j$, let $q=\#\{j:\alpha_j>0\}$, and write
\[
 \prod_{j=1}^{N}(s)_{\alpha_j}
      =\sum_{\ell=q}^{d}c_{\boldsymbol\alpha,\ell}s^\ell,
 \qquad (s)_m=s(s+1)\cdots(s+m-1),\quad (s)_0=1.
\]
Then
\begin{equation}
 \partial_{\boldsymbol\mu}^{\boldsymbol\alpha}J_{N,r}
 =(-1)^d r!\sum_{\ell=q}^{\min(d,r)}
       \frac{(-1)^{r-\ell}c_{\boldsymbol\alpha,\ell}}{(r-\ell)!}
       \sum_{n=0}^{\infty}
       \frac{S_n^{r-\ell}}
            {\prod_{j=1}^{N}(x_n^2+\mu_j)^{\alpha_j}}.
 \label{sp:all-mass-derivatives}
\end{equation}
An empty outer sum is zero.  Indeed, differentiating the spectral
summand gives
\[
 (-1)^d\prod_j(s)_{\alpha_j}\,e^{-sS_n}
                  \prod_j(x_n^2+\mu_j)^{-\alpha_j},
\]
whose series is normal near
zero; coefficient extraction gives
\eqref{sp:all-mass-derivatives}.  Each inner series is ordinary
and absolutely convergent.
\end{remark}

\begin{corollary}[Integral shifts and the zero-mass boundary]
\label{sp:integer-a}
For an integer $M\geq1$ and $t>0$,
\begin{equation}
 Q_M(t)=\frac{\pi\sqrt t\,\coth(\pi\sqrt t)-1}{2t}
           -\sum_{j=1}^{M-1}\frac{1}{j^2+t}.
 \label{sp:integer-resolvent}
\end{equation}
For integers $p\geq1$ its confluent derivatives give
\begin{equation}
 \sum_{n=0}^{\infty}\frac{1}{((n+M)^2+t)^p}
 =\frac{(-1)^{p-1}}{(p-1)!}\partial_t^{p-1}Q_M(t).
 \label{sp:confluent-resolvent}
\end{equation}
These expressions extend continuously to $t=0$, where
\begin{equation}
 \sum_{n=0}^{\infty}(n+M)^{-2p}
    =\zeta(2p)-H_{M-1}^{(2p)},
 \qquad H_m^{(v)}:=\sum_{j=1}^{m}j^{-v},\quad H_0^{(v)}:=0.
 \label{sp:harmonic-resolvent}
\end{equation}
Consequently, when $r\leq N$, the limit of
\eqref{sp:rational-mass-resolvent} as its $r$ selected masses tend
to zero is
\begin{equation}
 (-1)^r r!\{\zeta(2r)-H_{M-1}^{(2r)}\},\qquad a=M.
 \label{sp:harmonic-mixed-jet}
\end{equation}
The remaining masses do not enter this mixed derivative.
\end{corollary}

\begin{proof}
The partial-fraction expansion of the hyperbolic cotangent gives
\[
 \sum_{j=1}^{\infty}(j^2+t)^{-1}
       =\frac{\pi\sqrt t\coth(\pi\sqrt t)-1}{2t}.
\]
Removing its first $M-1$ terms proves
\eqref{sp:integer-resolvent}.  Normal convergence permits the
$t$ derivatives in \eqref{sp:confluent-resolvent}.  At $t=0$ the
series becomes a tail of the defining series for $\zeta(2p)$,
which proves \eqref{sp:harmonic-resolvent}.  Dominated convergence
in \eqref{sp:rational-mass-resolvent} gives the last assertion.
\end{proof}

\subsection{Compact densities for connected logarithmic products}
\label{sp:connected-subsection}

We now use $r$ for the number of logarithmic factors.  Fix an
anchor $\beta\geq0$, let $\nu_1,\ldots,\nu_r>\beta$, and put
$M=\max_j\nu_j$.  Define
\begin{equation}
 G_{\boldsymbol\nu,\beta}(z)
 :=\prod_{j=1}^{r}
       \log\frac{z+\nu_j}{z+\beta}.
 \label{sp:connected-product}
\end{equation}
Each logarithm means the analytic continuation from positive
real $z$ of $\Log(z+\nu_j)-\Log(z+\beta)$, with value tending
to zero at infinity.  This gives a holomorphic function on
$\mathbb C\setminus[-M,-\beta]$.
For $\beta<t<M$ away from the masses, write
\begin{equation}
 \ell_j(t):=\log\left|\frac{\nu_j-t}{\beta-t}\right|,
 \qquad \chi_j(t):=\boldsymbol1_{\{t<\nu_j\}},
 \qquad
 \rho_{\boldsymbol\nu,\beta}(t)
 :=-\frac1\pi\Im\prod_{j=1}^{r}
                     \{\ell_j(t)-i\pi\chi_j(t)\}.
 \label{sp:connected-density}
\end{equation}
The density is taken to be zero outside $[\beta,M]$.

\begin{theorem}[Compact representation and exact moments]
\label{sp:compact-density}
The density in \eqref{sp:connected-density} is integrable and
satisfies
\begin{equation}
 G_{\boldsymbol\nu,\beta}(z)
 =\int_\beta^M\frac{\rho_{\boldsymbol\nu,\beta}(t)}{z+t}\,dt,
 \qquad z\in\mathbb C\setminus[-M,-\beta].
 \label{sp:compact-stieltjes}
\end{equation}
For every integer $k\geq0$,
\begin{equation}
 \int_\beta^M t^k\rho_{\boldsymbol\nu,\beta}(t)\,dt
 =(-1)^k[u^{k+1}]
       \prod_{j=1}^{r}\log\frac{1+\nu_j u}{1+\beta u},
 \label{sp:compact-moments}
\end{equation}
where $[u^m]$ extracts a Taylor coefficient at $u=0$.
In particular, with $C:=\prod_{j=1}^{r}(\nu_j-\beta)$,
\begin{align}
 \int_\beta^M t^k\rho_{\boldsymbol\nu,\beta}(t)\,dt&=0,
       &&0\leq k\leq r-2,\label{sp:vanishing-moments}\\
 \int_\beta^M t^{r-1}\rho_{\boldsymbol\nu,\beta}(t)\,dt
      &=(-1)^{r-1}C.
 \label{sp:first-nonzero-moment}
\end{align}
The first range is empty if $r=1$.
\end{theorem}

\begin{proof}
On the upper cut $z=-t+i0$, the denominator logarithm has
imaginary part $\pi$ for $t>\beta$, while the numerator
logarithm has imaginary part $\pi$ precisely for $t>\nu_j$.
Their difference is therefore
$\ell_j(t)-i\pi\chi_j(t)$.
For $t>M$ all these differences have identical real boundary
values on both banks of the cut, which explains the compact
support.  At every endpoint the product has at worst a power of
a logarithm.  Thus the density is integrable, even when several
masses coincide, and every small endpoint circle contributes
$O(\varepsilon(1+|\log\varepsilon|^r))$ to the Cauchy integral.
At infinity,
\begin{equation}
 G_{\boldsymbol\nu,\beta}(z)=Cz^{-r}+O(z^{-r-1}).
 \label{sp:connected-asymptotic}
\end{equation}
The outer contour therefore vanishes.  Collapsing the Cauchy
contour onto the finite cut proves
\eqref{sp:compact-stieltjes}, with its indicated sign.

For $|z|>M$, expand $(z+t)^{-1}$ into its uniformly convergent
geometric series on the support of the density.  The coefficient
of $z^{-k-1}$ in the resulting integral is
$(-1)^k\int t^k\rho_{\boldsymbol\nu,\beta}(t)\,dt$.
On the other hand, the same coefficient on the left is
$[u^{k+1}]\prod_j\log((1+\nu_j u)/(1+\beta u))$.
This proves \eqref{sp:compact-moments}.  Each logarithmic factor
starts with $(\nu_j-\beta)u$, so the product starts with $Cu^r$.
The final two assertions follow.
\end{proof}

\begin{corollary}[Coincident masses]
\label{sp:coincident-density}
If $\nu_1=\cdots=\nu_r=\nu>\beta$, set
$y=\log((\nu-t)/(t-\beta))$ for $\beta<t<\nu$.  Then
\begin{equation}
 \rho_{\nu,\beta}^{(r)}(t)=P_r(y),\qquad
 P_r(y):=\sum_{h=0}^{\lfloor(r-1)/2\rfloor}
       (-1)^h\binom{r}{2h+1}\pi^{2h}y^{r-2h-1}.
 \label{sp:coincident-polynomial}
\end{equation}
For example,
\begin{equation}
 P_1(y)=1,\qquad P_2(y)=2y,\qquad
 P_3(y)=3y^2-\pi^2,\qquad P_4(y)=4y^3-4\pi^2y.
 \label{sp:coincident-examples}
\end{equation}
\end{corollary}

\begin{proof}
On this single interval every factor in
\eqref{sp:connected-density} equals $y-i\pi$.
Expanding $-\pi^{-1}\Im(y-i\pi)^r$ gives the formula.
\end{proof}

\begin{corollary}[An ordinary connected spectral sum]
\label{sp:ordinary-connected}
For $a>0$,
\begin{equation}
 \sum_{n=0}^{\infty}
   \prod_{j=1}^{r}\log\frac{x_n^2+\nu_j}{x_n^2+\beta}
 =\int_\beta^M\rho_{\boldsymbol\nu,\beta}(t)Q_a(t)\,dt.
 \label{sp:connected-sum}
\end{equation}
Both sides are ordinary and absolutely convergent.
\end{corollary}

\begin{proof}
The summand is $O(x_n^{-2r})$ by
\eqref{sp:connected-asymptotic}.  Moreover, $Q_a$ is bounded on
the compact interval $[\beta,M]$, including when $\beta=0$.
Thus $\int|\rho_{\boldsymbol\nu,\beta}(t)|Q_a(t)\,dt$ is finite,
and absolute Fubini applied to
\eqref{sp:compact-stieltjes} proves the assertion.
\end{proof}

\begin{remark}[A fixed number of spectral factors]
\label{sp:fixed-N-differences}
Suppose $N\geq r$ is fixed, and choose $r$ labeled mass slots.
In slot $j$, let $\Delta_j$ mean the value with mass $\nu_j$
minus the value with mass $\beta$; the other slots are held
fixed.  The normally convergent formula
\eqref{sp:normal-continuation} also permits zero masses at this
fixed positive $a$.  Consequently,
\begin{equation}
 \Delta_1\cdots\Delta_r J_{N,r}
 =(-1)^r r!\int_\beta^M
                    \rho_{\boldsymbol\nu,\beta}(t)Q_a(t)\,dt.
 \label{sp:connected-finite-difference}
\end{equation}
Indeed, in \eqref{sp:jet-ordinary-series} the mass-independent
baseline cancels, while the mixed difference of the $r$th power
of a sum of logarithms is $r!$ times the product of their $r$
increments.  This is a replacement operation at fixed $N$.
Adding or deleting a spectral factor changes the scale $2N$ in
the zeta term and is a different operation.
\end{remark}

\begin{theorem}[The critical odd weight as an ordinary sum]
\label{sp:critical-weight}
Under the assumptions of Theorem~\ref{sp:compact-density},
\begin{equation}
 \boxed{\begin{aligned}
 &\sum_{n=0}^{\infty}\left\{
        x_n^{2r-1}G_{\boldsymbol\nu,\beta}(x_n^2)-\frac{C}{x_n}
                         \right\}\\
 &\qquad=(-1)^r\int_\beta^M
       \rho_{\boldsymbol\nu,\beta}(t)t^{r-1}
       \{\Re\psi(a+i\sqrt t)-\psi(a)\}\,dt.
 \end{aligned}}
 \label{sp:critical-ordinary-sum}
\end{equation}
The summand on the left is $O(x_n^{-3})$.  The integral on the
right is an ordinary absolutely convergent integral; the
digamma difference has its continuous value zero at $t=0$.
\end{theorem}

\begin{proof}
Polynomial division gives
\[
 \frac{x^{2r-1}}{x^2+t}
 =\sum_{k=0}^{r-2}(-t)^k x^{2r-3-2k}
          +(-t)^{r-1}\frac{x}{x^2+t},
\]
where the sum is empty for $r=1$.
Upon integration against $\rho_{\boldsymbol\nu,\beta}$, every
polynomial term vanishes by \eqref{sp:vanishing-moments}.
Using \eqref{sp:first-nonzero-moment} for the subtraction gives
the exact identity, for each $x>0$,
\[
 x^{2r-1}G_{\boldsymbol\nu,\beta}(x^2)-\frac{C}{x}
 =(-1)^{r-1}\int_\beta^M
       \rho_{\boldsymbol\nu,\beta}(t)t^{r-1}
             \left\{\frac{x}{x^2+t}-\frac1x\right\}\,dt.
\]
The expression in braces equals $-t/[x(x^2+t)]$ and is bounded
in absolute value by $M/x^3$ on the whole integration interval.
Hence absolute Fubini is valid after this subtraction.
The convergent digamma difference series yields
\[
 \sum_{n=0}^{\infty}
       \left\{\frac{x_n}{x_n^2+t}-\frac1{x_n}\right\}
       =\psi(a)-\Re\psi(a+i\sqrt t).
\]
Substitution proves \eqref{sp:critical-ordinary-sum}.
The asserted $O(x_n^{-3})$ remainder also follows immediately
from \eqref{sp:connected-asymptotic}.
\end{proof}

\subsection{Polynomial weights and the inherited local correction}
\label{sp:weighted-bridge}

Let $W(x)=\sum_{d=0}^{D}w_dx^d$ be a polynomial.  For functions
with the relevant meromorphic continuation, use the fixed
regulator
\begin{equation}
 \mathcal R_W[f]
 :=\operatorname{FP}_{\varepsilon=0}
     \sum_{n=0}^{\infty}W(x_n)x_n^{-2\varepsilon}f(x_n),
 \qquad
 Q_{W,a}(t):=\mathcal R_W\left[x\longmapsto\frac{1}{x^2+t}\right],
 \label{sp:fixed-regulator}
\end{equation}
where $\operatorname{FP}$ means the constant Laurent coefficient.
The weight is already included in $\mathcal R_W$.

\begin{proposition}[The compact transform with a polynomial weight]
\label{sp:weighted-compact}
For the connected product in \eqref{sp:connected-product},
\begin{equation}
 \mathcal R_W[x\longmapsto G_{\boldsymbol\nu,\beta}(x^2)]
 =\int_\beta^M\rho_{\boldsymbol\nu,\beta}(t)Q_{W,a}(t)\,dt.
 \label{sp:weighted-compact-transform}
\end{equation}
The integral on the right is ordinary and absolutely convergent.
The kernels for monomials have the explicit expressions, for
integers $p\geq0$,
\begin{align}
 Q_{x^{2p},a}(t)
 &=\sum_{j=0}^{p-1}(-t)^j\zeta(2j-2p+2,a)
                        +(-t)^pQ_a(t),
 \label{sp:even-weight-kernel}\\
 Q_{x^{2p+1},a}(t)
 &=\sum_{j=0}^{p-1}(-t)^j\zeta(2j-2p+1,a)
                        -(-t)^p\Re\psi(a+i\sqrt t).
 \label{sp:odd-weight-kernel}
\end{align}
Both sums are empty for $p=0$, and the kernels have their
continuous values at $t=0$.
\end{proposition}

\begin{proof}
For $\Re\varepsilon>(D-1)/2$, the compact representation
\eqref{sp:compact-stieltjes} may be summed with the weight
$W(x_n)x_n^{-2\varepsilon}$ by absolute Fubini.
To continue this identity toward $\varepsilon=0$, choose $K$
sufficiently large and use the exact finite expansion
\[
 \frac{1}{x^2+t}
 =\sum_{j=0}^{K-1}(-t)^j x^{-2j-2}
       +\frac{(-t)^K}{x^{2K}(x^2+t)}.
\]
After summation the finite sum consists of terms
$w_d(-t)^j\zeta(2\varepsilon+2j+2-d,a)$.
The remainder is normally convergent for $\varepsilon$ in a
neighborhood of zero, uniformly on $\beta\leq t\leq M$.
Indeed, its summands have size at most a constant times
$x_n^{D-2\Re\varepsilon-2K-2}$ there.
The possible polar coefficients are polynomials in $t$, and
the holomorphic remainder has a uniform integrable majorant
against $|\rho_{\boldsymbol\nu,\beta}(t)|\,dt$.
Laurent coefficient extraction may therefore be interchanged
with the compact integral, proving
\eqref{sp:weighted-compact-transform}.  This argument uses
meromorphic continuation with uniform control, rather than a
pointwise limit of a divergent unregulated sum.

Polynomial division proves \eqref{sp:even-weight-kernel} from
$Q_{1,a}=Q_a$.  For the odd case, first observe that
\[
 \sum_{n=0}^{\infty}x_n^{-2\varepsilon}
                \frac{x_n}{x_n^2+t}
 =\zeta(1+2\varepsilon,a)
  +\sum_{n=0}^{\infty}x_n^{-2\varepsilon}
        \left\{\frac{x_n}{x_n^2+t}-\frac1{x_n}\right\}.
\]
The second series is normal near zero.  Since the constant
Laurent coefficient of $\zeta(1+2\varepsilon,a)$ is $-\psi(a)$,
the digamma difference used in
Theorem~\ref{sp:critical-weight} gives
$Q_{x,a}(t)=-\Re\psi(a+i\sqrt t)$.
Polynomial division now proves \eqref{sp:odd-weight-kernel}.
\end{proof}

\begin{remark}[Scope and normalization of the local term]
\label{sp:inherited-local-correction}
The ordinary scalar formula \eqref{sp:master-formula} and the
compact transform fit the local-correction calculation already
present in the arithmetic spectral jets continuation
\cite{ProveItArithmeticJets}.  We record its normalization here
as inherited input.  Define
\[
 J_{W,N,r}:=\left.\partial_s^r
       \sum_{n=0}^{\infty}W(x_n)
                     \prod_{j=1}^{N}(x_n^2+\mu_j)^{-s}
                    \right|_{s=0},\qquad
 L(u):=\sum_{j=1}^{N}\log(1+\mu_j u),
\]
with the spectral sum continued before taking its derivative.
For $r\geq1$, the existing correction formula is
\begin{equation}
 \begin{split}
 J_{W,N,r}
 &-\mathcal R_W\left[
        x\longmapsto\left(-\sum_{j=1}^{N}\log(x^2+\mu_j)\right)^r
                   \right]\\
 &=\frac{(-1)^{r+1}}{2N(r+1)}
     \sum_{k\geq r+1}w_{2k-1}[u^k]L(u)^{r+1}.
 \end{split}
 \label{sp:inherited-correction}
\end{equation}
The sum is finite because $W$ is a polynomial.  Its normalization
can be read from the resonant factor
$[u^k]e^{-sL(u)}\zeta(1+2Ns,a)$: the pole contributes
$r![s^{r+1}u^k]e^{-sL(u)}/(2N)$ to the $r$th spectral jet.
This equals the corresponding summand in
\eqref{sp:inherited-correction}.  The fixed regulator has no
such additional constant coefficient from its pole.

For $W=1$ the correction vanishes, consistent with the direct
normal-convergence proof of Theorem~\ref{sp:master}.
More generally, it vanishes for every even polynomial weight.
These are specializations of the inherited formula, not a new
general regularized-trace principle; all-coefficient Laurent
expansions of regularized traces belong to the established
framework represented by Paycha--Scott \cite{PaychaScott}.
Throughout such comparisons, interaction among masses means
replacement of labeled masses at fixed $N$, as in
Remark~\ref{sp:fixed-N-differences}.
\end{remark}

